\uuid{ka69}
\titre{Addition/multiplication}
\auteur{bodin ; hanani ; mzari}
\organisation{exo7}
\licence{CC BY-NC-SA 4.0}
\datecreate{2020-01-20}
\niveau{L1}
\module{Analyse}
\chapitre{Propriétés de R}
\sousChapitre{Propriétés des nombres réels}
\difficulte{1}
\theme{Propriétés de nombres réels, addition, multiplication}
\exo7qcmid{lille-1.reels.007}

\contenu{
\begin{question}
\qclassification{120.03}
\qtags{motcle=addition/multiplication}
Comment s'appelle les propriétés suivantes de $\Rr$ ?
\begin{answers}
    \good{$(a+b)+c=a+(b+c)$ est l'associativité de l'addition.}

    \bad{$(a \times b) \times c=a \times (b \times c)$ est la distributivité de la multiplication.}

    \good{$a \times b = b\times a$ est la commutativité de la multiplication.}

    \bad{$a \times (b+c) = a\times b + a \times c$ est l'intégrité.}
\end{answers}
\begin{explanations}
$(a+b)+c=a+(b+c)$ est l'associativité de l'addition.

$(a \times b)\times c=a \times (b \times c)$ est l'associativité de la multiplication.

$a \times b = b\times a$ est la commutativité de la multiplication.

$a \times (b+c) = a\times b + a \times c$ est la distributivité de la multiplication  par rapport à l'addition.
\end{explanations}
\end{question}
}
